Now we move from only 6 periods to 100 and 500 periods to determine where inertia is no longer dominant, and we just have the effect of the periodic input.
In order to keep computational times reasonable, the step-size is decreased from period/30 to period/10, i.e. the sinusoid input is divided into 30 and 10 segments, respectively.

Similar to before, figure \ref{fig:max-displacement-100} shows the maximum displacement over 100 periods.
In this format it is indistinguishable from figure \ref{fig:max-displacement-6} showing 6 periods, with the exception of around 900 Hz, as it approaches resonant frequency.
\begin{figure}[h]\centering
    \begin{tikzpicture}
        \begin{loglogaxis}[
            % width = 14cm, height = 6cm,
            xlabel = {Frequency [Hz]}, ylabel = {Maximum displacement [m]},
            title = {},
        ]
            \addplot+[only marks] table [ col sep=comma,]
                {Modelling/Frequencies-100-periods/spigot_maxima.csv};
            \addlegendentry{Spigot}
            \addplot+[ only marks ] table [col sep=comma,]
                {Modelling/Frequencies-100-periods/stem_maxima.csv};
            \addlegendentry{Stem}
            \addplot+[ only marks ] table [col sep=comma, ]
                {Modelling/Frequencies-100-periods/collar_maxima.csv};
            \addlegendentry{Collar}
        \end{loglogaxis}
    \end{tikzpicture}

\caption{Maximum displacement as a function of frequency over 100 periods}
\label{fig:max-displacement-100}
\end{figure}

Figure \ref{fig:long-term-displacement} (a) shows the first 500 periods of the stem for an input at 750 Hz.
(b) zooms into the first 100 periods to provide a better visualisation of what is happening.
We can clearly observe the damping over time, but it is hard to establish when the ripples are no longer significant.
In order to write software that quantifies the change over time, time is split into percentages beginning from the end and moving backwards, which is being called ``Tail-end percentage''.
E.g. The 10\% tail-end is the last 10\% of time, i.e. what happens between 90\% and 100\%.


\begin{figure}[h] \centering
    \begin{subfigure}{\textwidth}
        \begin{tikzpicture}[scale=1]
            \begin{axis}[
                width = 14cm, height = 6cm,
                xlabel = {Time [ms]}, ylabel = {Displacement [nm]},
            ]
                \addplot+[no marks] table [col sep=comma, x expr = 1e3* \thisrowno{0}, y expr = 1e9 * \thisrowno{3}]
                {Modelling/Frequencies-500-periods/csv/Job-19_750-Hz-500-periods_Node STEM-1.17.csv};
                \addlegendentry{Stem}
                \draw[red] (axis cs:600,-0.8) rectangle (axis cs:680,1.2);
                \addplot+[no marks] coordinates {(600,0)};
                \addlegendentry{Tail-end 10\%}
            \end{axis}
        \end{tikzpicture}
    \caption{First 500 periods, and highlighted }
    \end{subfigure}
    \begin{subfigure}{\textwidth}
            \begin{tikzpicture}[scale=1]
                \begin{axis}[
                    width = 14cm, height = 6cm,
                    xlabel = {Time [ms]}, ylabel = {Displacement [nm]},
                    ]
                    \addplot+[no marks] table [col sep=comma, x expr = 1e3* \thisrowno{0}, y expr = 1e9 * \thisrowno{3}]
                    {Modelling/Frequencies-100-periods/csv/Job-42_750-Hz-100-periods_Node STEM-1.17.csv};
                    \addlegendentry{Stem}
                \end{axis}
            \end{tikzpicture}
            \caption{First 100 periods}
        \end{subfigure}
\caption{Displacement as a function of time for sinusoidal input at \SI{750}{Hz}. These images highlight damping over the time, though determining when they become insignificant is challenging.}
\label{fig:long-term-displacement}
\end{figure}


Finding the maxima across all intervals allow us to ascertain 1) the largest and smallest maxima at each frequency, and 2) determine the time-invariant frequencies.
Figure \ref{fig:long-term-rel-diff} shows precisely this, with one peculiarity: Because we are approaching resonance at \SI{900}{Hz}, the maxima from tail-end is always the global maximum.
This is shown in the graph as 0\% change at \SI{900}{Hz}, but it is time-variant at that frequency.
In any case, we do not want an ITAP vibrating close to resonance, so we can focus on quantifying the change over time for the well behaved time-variant frequencies.

Knowing the time-variant frequencies, we can use the tail-end subdivisions to find the maximum displacements for all intervals, which we compare to the amplitude at the final periods where there are no ripples.
The result of this is quantifying the relative loss of amplitude over time, which is seen in Figure \ref{fig:long-term-abs-diff}.
While the time-variant frequencies are perhaps outside the frequencies of clinical use, this characterisation combined with the frequency response gives great information about the amplitude change over time.
To the point that it might be worth exploring higher frequencies if more granular control of the amplitude becomes necessary in the future. 

\begin{figure}[h] \centering
    \begin{tikzpicture}[scale=1]
        \begin{axis}[
            width = 14cm, height = 6cm,
            % title = {Relative difference in maximum displacement before and after damping},
            xlabel = {Frequency [Hz]}, ylabel = {Relative amplitude decrease [\%]},
            ymax = 100,
            ytick = {0, 20,...,100},
            % grid = both, minor tick num = 5,
        ]
        \addplot+[no marks] table [col sep=comma, y expr = 100*\thisrowno{1}]
            {Modelling/Error/spigot_rel_diff.csv};
        \addlegendentry{Spigot}
        \addplot+[no marks] table [col sep=comma, y expr = 100*\thisrowno{1}]
            {Modelling/Error/stem_rel_diff.csv};
        \addlegendentry{Stem}
        \addplot+[no marks] table [col sep=comma, y expr = 100*\thisrowno{1}]
            {Modelling/Error/collar_rel_diff.csv};
        \addlegendentry{Collar}
        \end{axis}
    \end{tikzpicture}
\caption{Difference between largest and smallest maxima observed throughout 500-periods. The simulation is partitioned by percentage of time elapsed (starting from the end) }
\label{fig:long-term-rel-diff}
\end{figure}

\begin{figure}[h] \centering
    \begin{subfigure}{0.32\textwidth}
        \begin{tikzpicture}[scale=0.65]
            \begin{axis}[
                % width = 14cm, height = 6cm,
                xlabel = {Percentage of time elapsed}, ylabel = {Relative amplitude loss },
                legend columns = 1,
                legend cell align = left,
                legend pos = outer north east,
                ymax = 35,
            ]
            \foreach \freq in {450,500,...,850} {
                \addplot+[no marks] table [col sep=comma, x expr = 100*(1-\thisrowno{0}), y expr = 100*\thisrowno{1}] {Modelling/Error/spigot-\freq.csv};
                % \expandafter\addlegendentry\expandafter{\freq Hz}
            };
            \end{axis}
        \end{tikzpicture}
        \caption{Spigot}
    \end{subfigure}%
    \begin{subfigure}{0.3\textwidth}
        \begin{tikzpicture}[scale=0.65]
            \begin{axis}[
                % width = 14cm, height = 6cm,
                yticklabels = \empty,
                xlabel = {Percentage of time elapsed}, %ylabel = {Relative amplitude loss },
                legend columns = 1,
                legend cell align = left,
                legend pos = outer north east,
                ymax = 35,
            ]
            \foreach \freq in {450,500,...,850} {
                \addplot+[no marks] table [col sep=comma, x expr = 100*(1-\thisrowno{0}), y expr = 100*\thisrowno{1}] {Modelling/Error/stem-\freq.csv};
                % \expandafter\addlegendentry\expandafter{\freq Hz}
            };
            \end{axis}
        \end{tikzpicture}
        \caption{Stem}
    \end{subfigure}%
    \begin{subfigure}{0.37\textwidth}
        \begin{tikzpicture}[scale=0.65]
            \begin{axis}[
                % width = 14cm, height = 6cm,
                yticklabels = \empty,
                xlabel = {Percentage of time elapsed}, %ylabel = {Relative amplitude loss },
                legend columns = 1,
                legend cell align = left,
                legend pos = outer north east,
                ymax = 35,
            ]
            \foreach \freq in {450,500,...,850} {
                \addplot+[no marks] table [col sep=comma, x expr = 100*(1-\thisrowno{0}), y expr = 100*\thisrowno{1}] {Modelling/Error/collar-\freq.csv};
                \expandafter\addlegendentry\expandafter{\freq Hz}
            };
            \end{axis}
        \end{tikzpicture}
        \caption{Collar}
    \end{subfigure}
    \caption{Amplitude decline over time for time-variant frequencies. By picking a particular frequency, we can compare (a), (b) and (c) to }
\label{fig:long-term-abs-diff}
\end{figure}